\begin{exercisechunk}
\item \textbf{Denoising coefficients by parameter sharing.}
\label[exercise]{ex:l5-sharing}
Let \(U\in\R^{p\times r}\) be known with \(U^\top U=I_r\).
Suppose a preliminary estimate is
\(\widehat b=b+\xi\), where \(\xi\sim\mathcal N(0,\sigma^2I_p/n)\).
Set \(\widetilde b=UU^\top\widehat b\).
\begin{enumerate}[(a),beginpenalty=10000]
\item If \(b\) lies in the column space of \(U\), compute the expected
squared coefficient errors of \(\widehat b\) and \(\widetilde b\).
\item Let an independent test feature \(\Psi\) have mean zero and
\(\E[\Psi\Psi^\top]=I_p\). Compute the corresponding prediction
errors for predicting \(b^\top\Psi\).
\item Realize a one-dimensional example using
\(b_j=\theta\lambda^{j-1}\), \(j=1,\ldots,p\), with known
\(|\lambda|<1\). Explain the relationship to a scalar recurrent predictor.
\item For arbitrary \(b\), derive the bias--variance formula for
\(\widetilde b\). Exactly when is its expected coefficient error
smaller than that of \(\widehat b\)?
\end{enumerate}
The independent isotropic coefficient errors are a simplifying model,
not an assertion about lagged time-series least squares.
\end{exercisechunk}
